\section{总结与延伸}

Lecture 8 是 CS336 中把分布式训练从“API 原语”推到“系统设计”的关键一讲。它把 Lecture 7 的 collectives 放进真实 LLM 训练：DDP 用 all-reduce；ZeRO/FSDP 用 reduce-scatter 和 all-gather；pipeline 用 send/recv 或 stage 边界通信；tensor parallel 在每层通信；sequence/context parallel 处理 activation 和长上下文；expert parallel 则围绕 all-to-all 和负载均衡展开。

\begin{importantbox}{最终 takeaways}
\begin{enumerate}
    \item 单 GPU scaling 同时受 compute 和 memory 限制，训练单位会自然扩大到 datacenter。
    \item collectives 是并行训练的底层语言；all-reduce = reduce-scatter + all-gather 是理解 ZeRO/FSDP 的核心。
    \item ZeRO-1/2 主要降低 data parallel 的状态重复，ZeRO-3/FSDP 进一步分片参数，但通信更进入关键路径。
    \item pipeline parallel 沿 depth 切，省参数显存且可跨慢链路，但要处理 bubble。
    \item tensor parallel 沿 width 切，利用率高但通信高频，通常限制在高速节点内。
    \item sequence/context/expert parallelism 分别处理 activation、长上下文和 MoE experts。
    \item 真实训练常用 3D/4D 组合；配置的正确性来自状态和通信账本，而不是并行名词数量。
\end{enumerate}
\end{importantbox}

\subsection{拓展阅读}

阅读并行训练材料时，建议为每个方案建立同一张账本：长期驻留的 parameters/gradients/optimizer states/activations，各通信 primitive 的频率与 bytes，通信所在硬件域，以及 global/micro batch 的约束。这样才能把不同论文中的配置还原成可比较的系统选择。

\begin{longtable}{L{0.23\textwidth}L{0.31\textwidth}L{0.36\textwidth}}
\toprule
材料 & 为什么值得读 & 带着什么问题读 \\
\midrule
CS336 Lecture 7--8 slides & 从 collective primitives 过渡到完整并行配置 & 每种并行最终分解成哪些 all-reduce、all-gather、reduce-scatter、all-to-all？ \\
Rajbhandari et al., ZeRO & 建立 optimizer/gradient/parameter state sharding 的原始内存账本 & ZeRO-1/2 为什么近似不增加总通信，ZeRO-3 为什么进入关键路径？ \\
PyTorch FSDP tutorial & 展示参数 materialization、module wrapping、prefetch 与 reduce-scatter 的实际生命周期 & Wrap 粒度怎样影响峰值 buffer、消息数量和 overlap？ \\
Megatron-LM / Megatron-Core docs & 覆盖 TP、PP、SP、CP、EP 与 MoE 并行 group 的工程实现 & 哪些通信必须留在节点内，attention 与 experts 如何使用不同 group？ \\
Narayanan et al. 2021 & 提供 3D parallel scaling、pipeline schedules 与 TP=8 等经验证据 & 线性 scaling 在什么 batch、拓扑和重计算条件下成立？ \\
公开模型技术报告 & DeepSeek、Llama 3、Gemma、Mixtral、Nemotron、Qwen 等真实配置 & 并行数字是否自洽地乘到 world size，哪些字段可能未披露或随阶段变化？ \\
\bottomrule
\caption{Lecture 8 的注释式拓展阅读路线。}
\end{longtable}

\begin{importantbox}{建议的配置练习}
给定 64 台、每台 8×80GB GPU 的集群，分别为 70B dense、400B dense、200B-A20B MoE 与 1M-context MoE 设计初始 DP/TP/PP/EP/CP 配置。先证明状态能放入 HBM，再估计高频 collective 是否跨节点，最后用 microbatch、recomputation 和 overlap 调整吞吐。答案不要求唯一，但每个数字都必须由 memory、bandwidth 或 batch constraint 支持。
\end{importantbox}

\end{document}
